\documentclass[10pt,a4paper]{amsart}
\usepackage[margin=19mm]{geometry}
\usepackage{amsmath,amssymb,mathtools}
\usepackage[scr=rsfs,cal=euler]{mathalfa}
\usepackage{xfrac}
\usepackage[unicode,hidelinks,bookmarks=false]{hyperref}
\newcommand{\ie}{i.e.\ }
\newcommand{\cf}{cf.\ }
\newcommand{\resp}{respectively\ }
\newcommand{\Subsection}{\subsection}
\providecommand{\nobreakdash}{\nobreak\hbox{-}}
\setlength{\emergencystretch}{2em}
\multlinegap0pt
\usepackage{amssymb}
\usepackage{float}
\usepackage{enumerate}
\newtheorem{theorem}{Theorem}
\newtheorem{lemma}{Lemma}
\usepackage{cleveref}
\crefname{theorem}{Theorem}{Theorems}
\crefname{lemma}{Lemma}{Lemmas}
\newcommand{\N}{\mathbb{N}}
\newcommand{\R}{\mathbb{R}}
\renewcommand{\d}{\textnormal{d}}
\title{Local regularity for nonlocal equations with variable exponents}
\author{Jamil Chaker, Minhyun Kim}
\date{Source: arXiv:2107.06043v2 (CC BY 4.0)}
\begin{document}
\maketitle

\noindent\emph{Source and licence.} Local regularity for nonlocal equations with variable exponents. Version arXiv:2107.06043v2 (CC BY 4.0), distributed by arXiv under the Creative Commons Attribution 4.0 International licence, http://creativecommons.org/licenses/by/4.0/, as recorded in the arXiv licence metadata for this exact version and shown on the abstract page https://arxiv.org/abs/2107.06043v2. Full licence notice: \texttt{licenses/arxiv-2107.06043-CC-BY-4.0.txt}.\par\smallskip

\noindent\textit{Editorial scope.} Counted unit: the article's theorem thm:Holder (source0.tex lines 252--254). The article's complete original proof of it is delivered with it (source0.tex lines 1089--1199, one \texttt{proof} environment of 111 lines, which the article places away from the statement). No mathematics is written, repaired or re-derived: the statement and the proof are the article's own lines, in the article's own notation, and no retained line is substituted or re-flowed.  Retained uncounted as the source-local context the proof needs: lines 110--112; lines 137--139; lines 150--152; lines 236--242; lines 238--240; lines 849--851; lines 853--855; lines 949--956.  Nothing was omitted from any retained span, so the package carries no omission record.  No retained line contains a citation, so the wrapper carries no bibliography and no bibliography entry.  No retained line contains a figure environment, an image-inclusion command or drawing code, so no image is delivered and the package declares no asset dependency.  Editorial formatting only, in addition: an AMS \texttt{amsart} preamble that restates the article's own declarations and macros used by the retained lines, namely the environments \texttt{theorem}, \texttt{lemma} on separate counters, together with the standard packages those lines need.  The retained environments print in this document as Theorem 1, Lemma 1 in order of appearance, whereas the article prints the same items with the numbering of its own full document.  The complete native source of the article as distributed in the arXiv e-print for this exact version is bundled unchanged as \texttt{sources/113755/source0.tex}.
\begin{equation} \label{eq:EL}
(-\Delta)^s_{p(\cdot,\cdot)}u = 0
\end{equation}

\begin{equation} \label{eq:P1} \tag{P1}
R^{p_-(B \times B) - p_+(B \times B)} \leq L \quad \text{for all}~ \overline{B} = \overline{B_R(x_0)} \subset \Omega,
\end{equation}

\begin{equation} \label{eq:P2} \tag{P2}
p_+(B \times B^c) \leq p_+(B\times B) \quad\text{and}\quad p_-(B \times B^c) \leq p_-(B\times B) \quad\text{for all}~ \overline{B} = \overline{B_R(x_0)} \subset \Omega.
\end{equation}

\begin{theorem} \label{thm:local_boundedness}
Let $\Omega$ be a bounded domain in $\mathbb{R}^n$. If $u \in W^{s, p(\cdot,\cdot)}(\mathbb{R}^n)$ satisfies
\begin{equation} \label{eq:tail}
\sup_{x\in \Omega} \int_{\mathbb{R}^n} \frac{u_+(y)^{p(x,y)-1}}{1+|y|^{n+sp(x,y)}} \,\mathrm{d}y < +\infty \quad \left( \sup_{x\in \Omega} \int_{\mathbb{R}^n} \frac{u_-(y)^{p(x,y)-1}}{1+|y|^{n+sp(x,y)}} \,\mathrm{d}y < +\infty,~ \text{respectively} \right)
\end{equation}
and is a weak subsolution (supersolution, respectively) to \eqref{eq:EL} in $\Omega$, then $u$ is locally bounded from above (below, respectively) in $\Omega$. If $u \in W^{s, p(\cdot,\cdot)}(\mathbb{R}^n)$ is a weak solution to \eqref{eq:EL} in $\Omega$ satisfying \eqref{eq:tail} with $u_+$ replaced by $|u|$, then $u \in L^\infty_{\mathrm{loc}}(\Omega)$.
\end{theorem}

\begin{equation} 
\sup_{x\in \Omega} \int_{\mathbb{R}^n} \frac{u_+(y)^{p(x,y)-1}}{1+|y|^{n+sp(x,y)}} \,\mathrm{d}y < +\infty \quad \left( \sup_{x\in \Omega} \int_{\mathbb{R}^n} \frac{u_-(y)^{p(x,y)-1}}{1+|y|^{n+sp(x,y)}} \,\mathrm{d}y < +\infty,~ \text{respectively} \right)
\end{equation}

\begin{equation} \label{eq:assumption1}
0 \leq u \leq 2H \quad\text{in}~ B_R \quad\text{and}\quad |B_{R/2} \cap \lbrace u \geq H \rbrace| \geq \gamma |B_{R/2}|
\end{equation}

\begin{equation} \label{eq:assumption_tail2}
\sup_{x \in B_{3R/4}} \int_{\mathbb{R}^n \setminus B_R} \frac{u_-(y)^{p(x,y)-1}}{|y-x_0|^{n+sp(x,y)}} \,\mathrm{d}y \leq R^{-sp_+} (\delta H)^{p_+-1} + R^{-sp_-} (\delta H)^{p_--1}.
\end{equation}

 \begin{lemma} \label{lem:growth}
Let $B_R = B_R(x_0) \subset \mathbb{R}^n$ with $R \in (0,1)$. Let $H > 0$, and $0<\sigma < s <1$. Assume that $p$ satisfies \eqref{eq:P1} and \eqref{eq:P2} in $B_R$, $H^{p_+ - p_-} \leq 2$, and $p_+ < p_-^\ast$, where $p_\pm = p_\pm(B_R \times B_R)$ and $p_-^\ast = \frac{np_-}{n-\sigma p_-}$. Let $u \in W^{s, p(\cdot,\cdot)}(\mathbb{R}^n)$ be a weak supersolution to \eqref{eq:EL} in $B_R$ such that
\eqref{eq:assumption1} is satisfied for some $\gamma\in(0,1)$. Then there exists $\delta\in(0,1/8]$, such that, if $R^s\leq \delta H$ and \eqref{eq:assumption_tail2} is satisfied, then
\begin{equation}\label{eq:lemgrowth}
 u \geq \delta H \qquad \text{in } B_{R/4}.
 \end{equation}
 The constant $\delta$ depends on $n$, $s$, $\sigma$, $p_+(B_R \times \mathbb{R}^n)$, $p_-(B_R \times \mathbb{R}^n)$, $q$, and $L$.
\end{lemma}

\begin{theorem} \label{thm:Holder}
Let $\Omega$ be a bounded domain in $\mathbb{R}^n$. Assume that $p$ satisfies \eqref{eq:P1} and \eqref{eq:P2} in $\Omega$. If $u \in W^{s, p(\cdot,\cdot)}(\mathbb{R}^n)$ is a weak solution to \eqref{eq:EL} in $\Omega$ satisfying \eqref{eq:tail}, then $u$ is locally H\"older continuous in $\Omega$.
\end{theorem}

\begin{proof} [Proof of \Cref{thm:Holder}]
Let $x_0\in\R^n$. If $p(x_0, x_0) > n/s$, then we can find $R > 0$ and $\alpha \in (0,1)$ such that $B_R(x_0) \Subset \Omega$ and $u \in C^\alpha(\overline{B_{R}(x_0)})$ as in the proof of \Cref{thm:local_boundedness}. Thus, let us assume $p(x_0, x_0) \leq n/s$ in the rest of the proof. In this case, for given $\sigma \in (0,s)$ we can find $R \in (0,1)$ such that $B_R(x_0) \Subset \Omega$ and $p_+(B_R(x_0)\times B_R(x_0)) < p_-^{\ast}(B_R(x_0)\times B_R(x_0))$, where $p_-^{\ast}(B_R(x_0)\times B_R(x_0))= \frac{np_-(B_R(x_0)\times B_R(x_0))}{n-\sigma p_-(B_R(x_0)\times B_R(x_0))}$. By \Cref{thm:local_boundedness}, $u\in L^{\infty}(B_R(x_0))$.

Let $\delta\in(0,1)$ be the constant from \Cref{lem:growth} and let 
\begin{equation}\label{eq:alphabounds}
0<\alpha < \min\left\{ s,\log_4\left(\frac{2}{2-\delta}\right), \frac{sp_{+}(\Omega\times\R^n)}{2(p_{+}(\Omega\times\R^n)-1)}\right\}
\end{equation}
be chosen such that the following is satisfied:
\begin{equation}\label{eq:propalpha1}
\int^{\infty}_1 \frac{((4t)^{\alpha} - 1)^{p_{+}(\Omega\times \R^n)-1}}{t^{1+sp_{+}(\Omega\times\R^n)}}\, \d t + \int^{\infty}_1 \frac{((4t)^{\alpha} - 1)^{p_{-}(\Omega\times \R^n)-1}}{t^{1+sp_{-}(\Omega\times\R^n)}}\, \d t \leq \frac{\delta^{p_{+}(\Omega\times\R^n)-1}}{2^{p_{+}(\Omega\times\R^n)}n\omega_n},
\end{equation}
where $\omega_n$ denotes the volume the $n$-dimensional Euclidean unit ball.
We define $j_0\in\N$ to be the smallest natural number satisfying
\begin{equation}\label{eq:j0}
\begin{aligned}
j_0\geq \max\Bigg\{ &\frac{sp_{+}(\lbrace x_0 \rbrace \times B_R^c)}{2}\left|\log_4\left(\frac{\delta^{p_{+}(\lbrace x_0 \rbrace \times B_R^c)-1}}{2C_0}\right) \right|, \\
&\frac{sp_{-}(\Omega \times \R^n)}{2}\left|\log_4\left(\frac{\delta^{p_{-}(\Omega \times \R^n)-1}}{2C_0}\right) \Bigg|, \frac{|\log_4 (\frac{\delta}{2})|}{s-\alpha} \right\},
\end{aligned}
\end{equation}
where $C_0:=\max\{1,2^{p_{+}(\Omega\times\R^n)}\}\left(\frac{n\omega_n}{sp_{-}(\Omega\times\R^n)}+1\right)$.

In the following we show that there is a non-increasing sequence $(M_j)$ and a non-decreasing sequence $(m_j)$ in $\R$, such that for all $j\in\N\cup\{0\}$ 
\begin{equation}\label{eq:sequencesMjmj}
m_j \leq u \leq M_j \quad \text{in } B_{4^{-j}R}(x_0) \qquad \text{and} \qquad M_j - m_j = Z4^{-\alpha j},
\end{equation}
where 
\begin{align*}
Z&:= 2\cdot 4^{\alpha j_0} \|u\|_{L^{\infty}(B_R(x_0))} + R^s + 1 \\
& \quad + \left(R^{sp_{+}(\{x_0\}\times B_R(x_0)^c)} \sup_{x \in B_{3R/4}(x_0)} \int_{\mathbb{R}^n \setminus B_R(x_0)} \frac{|u(y)|^{p(x,y)-1}}{|y-x_0|^{n+sp(x,y)}} \,\mathrm{d}y\right)^{\frac{1}{p_{+}(\{x_0\}\times B_R(x_0)^c)-1}} . 
\end{align*}
For $j\in \{0,\dots,j_0\}$, we define $M_j:=4^{-\alpha j} Z/2$ and $m_j:=-4^{-\alpha j} Z/2$. Then \eqref{eq:sequencesMjmj} is clearly satisfied for all $j\in \{0,\dots,j_0\}$. It remains to prove the assertion \eqref{eq:sequencesMjmj} for $j > j_0$. 
The proof of the assertion follows by induction. Let us fix $j\geq j_0$ and assume that \eqref{eq:sequencesMjmj} is true for all $i\in \{0,\dots,j\}$. We now construct the elements $m_{j+1}$ and $M_{j+1}$ of the sequences. We distinguish between two cases.  

First, we assume 
\begin{equation}\label{eq:Hoeldcase1}
 \left| B_{\frac{4^{-j}R}{2}}(x_0) \cap \left\{ u\geq m_{j} + \frac{M_j-m_j}{2} \right\} \right| \geq \frac{1}{2}\left| B_{\frac{4^{-j}R}{2}}(x_0) \right|. 
\end{equation}
In this case, we define $v:=u-m_j$, $H:=\frac{M_j-m_j}{2}$ and $\widetilde{R}:=4^{-j}R$. \\
The main idea for constructing $m_{j+1}$ and $M_{j+1}$ is to apply \Cref{lem:growth} for the function $v$ and the radius $\widetilde{R}$. Hence, we need to verify the requirements of the lemma.
Note that by assumption we have $0\leq v\leq 2H$ in $B_{\widetilde{R}}(x_0)$ and  $| B_{\widetilde{R}/2}(x_0) \cap \left\{ v\geq H \right\} | \geq \frac{1}{2}| B_{\widetilde{R}/2}(x_0)|.$ It remains to prove $\widetilde{R}^s\leq \delta H$ and \eqref{eq:assumption_tail2}. First we show $\widetilde{R}^s\leq \delta H$. Note that $2H = M_j - m_j = Z4^{-\alpha j} \geq R^s 4^{-\alpha j}$. Since $j\geq j_0$, we can use \eqref{eq:j0}, which leads to
\[ \widetilde{R}^s = 4^{-js} R^s \leq  4^{j(\alpha-s)}2H \leq \delta H.\]
It remains to prove \eqref{eq:assumption_tail2}. We split $\R^n\setminus B_{\widetilde{R}}(x_0)$ as follows:
\[ \R^n\setminus B_{\widetilde{R}}(x_0) = \left(\R^n\setminus B_R(x_0)\right) \cup \left(\bigcup_{l=0}^{j-1} B_{4^{-l}R}(x_0) \setminus B_{4^{-(l+1)}R}(x_0)\right). \]
If $x\in B_{4^{-l}R}(x_0) \setminus B_{4^{-(l+1)}R}(x_0)$, then $|x-x_0|\geq 4^{-l-1}R$ and therefore
\begin{align*}
v(x) &= u(x)-m_j \geq m_l-M_l + 2H  = 2H(-4^{(-l+j)\alpha}+1) \\
& \geq -2H\left(\left( \frac{4|x-x_0|}{\widetilde{R}} \right)^{\alpha}-1\right). 
\end{align*}
On the other hand, if $x\in \R^n\setminus B_R(x_0)$, we have $v(x) \geq -|u(x)|-Z/2$.
 Now we are in a position to finalize the verification of \eqref{eq:assumption_tail2}. By the previous estimates on $v$ in $\R^n\setminus B_{\widetilde{R}}(x_0)$, we have
\begin{align*}
&\sup_{x \in B_{3\widetilde{R}/4}(x_0)} \int_{\mathbb{R}^n \setminus B_{\widetilde{R}}(x_0)} \frac{v_-(y)^{p(x,y)-1}}{|y-x_0|^{n+sp(x,y)}} \,\mathrm{d}y \\
&\leq \sup_{x \in B_{3\widetilde{R}/4}(x_0)} \int_{\R^n \setminus B_{\widetilde{R}}(x_0)} \frac{\left(2H\left(\left( \frac{4|y-x_0|}{\widetilde{R}} \right)^{\alpha}-1\right)\right)^{p(x,y)-1} }{|y-x_0|^{n+sp(x,y)}} \,\mathrm{d}y  \\
&\qquad +\max\{1,2^{p_+(B_R(x_0)\times\R^n)-1}\}\sup_{x \in B_{3\widetilde{R}/4}(x_0)} \int_{\mathbb{R}^n \setminus B_{R}(x_0)} \frac{|u(y)|^{p(x,y)-1}+Z^{p(x,y)-1}}{|y-x_0|^{n+sp(x,y)}} \,\mathrm{d}y\\
&  =:J_1+J_2.
\end{align*}
First note, that we can estimate $J_1$ as follows:
\begin{align*}
J_1&\leq \int_{\mathbb{R}^n \setminus B_{\widetilde{R}}(x_0)} \,\frac{\left(2H\left(\left( \frac{4|y-x_0|}{\widetilde{R}} \right)^{\alpha}-1\right)\right)^{p_+(B_{\widetilde{R}}(x_0) \times B_{\widetilde{R}}(x_0)^c)-1}}{|y-x_0|^{n+sp_+(B_{\widetilde{R}}(x_0) \times B_{\widetilde{R}}(x_0)^c)}} \, \d y\\
&\quad + \int_{\mathbb{R}^n \setminus B_{\widetilde{R}}(x_0)} \frac{\left(2H\left(\left( \frac{4|y-x_0|}{\widetilde{R}} \right)^{\alpha}-1\right)\right)^{p_-(B_{\widetilde{R}}(x_0) \times B_{\widetilde{R}}(x_0)^c)-1}}{|y-x_0|^{n+sp_-(B_{\widetilde{R}}(x_0) \times B_{\widetilde{R}}(x_0)^c)}}\, \d y\\
& = n\omega_n (2H)^{p_+(B_{\widetilde{R}}(x_0) \times B_{\widetilde{R}}(x_0)^c)-1} \widetilde{R}^{-sp_+(B_{\widetilde{R}}(x_0) \times B_{\widetilde{R}}(x_0)^c)}\int^{\infty}_1 \frac{((4t)^{\alpha} - 1)^{p_{+}(B_{\widetilde{R}}(x_0)\times B_{\widetilde{R}}(x_0)^c)-1}}{t^{1+sp_{+}(B_{\widetilde{R}}(x_0)\times B_{\widetilde{R}}(x_0)^c)}}\, \d t \\
& \quad + n\omega_n (2H)^{p_-(B_{\widetilde{R}}(x_0) \times B_{\widetilde{R}}(x_0)^c)-1} \widetilde{R}^{-sp_-(B_{\widetilde{R}}(x_0) \times B_{\widetilde{R}}(x_0)^c)}\int^{\infty}_1 \frac{((4t)^{\alpha} - 1)^{p_{-}(B_{\widetilde{R}}(x_0)\times B_{\widetilde{R}}(x_0)^c)-1}}{t^{1+sp_{-}(B_{\widetilde{R}}(x_0)\times B_{\widetilde{R}}(x_0)^c)}}\, \d t.
\end{align*}
Using $p_{-}(\Omega\times\R^n)\leq p_{\pm}(B_{\widetilde{R}}(x_0)\times B_{\widetilde{R}}(x_0)^c) \leq p_{+}(\Omega\times\R^n)$ and \eqref{eq:propalpha1}, we get 
\begin{align*}
\int^{\infty}_1& \frac{((4t)^{\alpha} - 1)^{p_{\pm}(B_{\widetilde{R}}(x_0)\times B_{\widetilde{R}}(x_0)^c)-1}}{t^{1+sp_{\pm}(B_{\widetilde{R}}(x_0)\times B_{\widetilde{R}}(x_0)^c)}}\, \d t \\
&\leq \int^{\infty}_1 \frac{((4t)^{\alpha} - 1)^{p_{+}(\Omega\times \R^n)-1}}{t^{1+sp_{+}(\Omega\times \R^n)}}\, \d t
+ \int^{\infty}_1 \frac{((4t)^{\alpha} - 1)^{p_{-}(\Omega\times \R^n)-1}}{t^{1+sp_{-}(\Omega\times \R^n)}}\, \d t \\
& \leq \frac{\delta^{p_{+}(\Omega\times\R^n)-1}}{2^{p_{+}(\Omega\times\R^n)}n\omega_n}.
\end{align*}
Combing the previous two estimates, we arrive at 
\begin{align*}
J_1 & \leq  \frac12 H^{p_+(B_{\widetilde{R}}(x_0) \times B_{\widetilde{R}}(x_0)^c)-1} \widetilde{R}^{-sp_+(B_{\widetilde{R}}(x_0) \times B_{\widetilde{R}}(x_0)^c)} \delta^{p_{+}(\Omega\times\R^n)-1} \\
& \quad +\frac12 H^{p_-(B_{\widetilde{R}}(x_0) \times B_{\widetilde{R}}(x_0)^c)-1} \widetilde{R}^{-sp_-(B_{\widetilde{R}}(x_0) \times B_{\widetilde{R}}(x_0)^c)}\delta^{p_{+}(\Omega\times\R^n)-1} \\
& \leq \frac12 (\delta H)^{p_+(B_{\widetilde{R}}(x_0) \times B_{\widetilde{R}}(x_0)^c)-1} \widetilde{R}^{-sp_+(B_{\widetilde{R}}(x_0) \times B_{\widetilde{R}}(x_0)^c)}  \\
& \quad +\frac12 (\delta H)^{p_-(B_{\widetilde{R}}(x_0) \times B_{\widetilde{R}}(x_0)^c)-1} \widetilde{R}^{-sp_-(B_{\widetilde{R}}(x_0) \times B_{\widetilde{R}}(x_0)^c)} \\
& \leq \frac12 (\delta H)^{p_+(B_{\widetilde{R}}(x_0) \times B_{\widetilde{R}}(x_0))-1} \widetilde{R}^{-sp_+(B_{\widetilde{R}}(x_0) \times B_{\widetilde{R}}(x_0))}  \\
& \quad +\frac12 (\delta H)^{p_-(B_{\widetilde{R}}(x_0) \times B_{\widetilde{R}}(x_0))-1} \widetilde{R}^{-sp_-(B_{\widetilde{R}}(x_0) \times B_{\widetilde{R}}(x_0))}.
\end{align*}
In the last inequality we used that by \eqref{eq:P2} and $\widetilde{R}^s\leq \delta H$,
\[  R^{-s p_{\pm}(B_R(x_0)\times B_R(x_0)^c)} (\delta H)^{p_{\pm}(B_R(x_0)\times B_R(x_0)^c)}   \leq R^{-s p_{\pm}(B_R(x_0)\times B_R(x_0))} (\delta H)^{p_{\pm}(B_R(x_0)\times B_R(x_0))}. \]
Next, we estimate $J_2$ as follows:
\begin{equation}\label{eq:estJ2}
\begin{aligned}
J_2 &\leq \max\{1,2^{p_+(\Omega\times\R^n)-1}\} R^{-sp_{+}(\{x_0\}\times B_R(x_0)^c)}Z^{p_{+}(\{x_0\}\times B_R(x_0)^c)-1} \\
&\quad + \max\{1,2^{p_+(\Omega\times\R^n)-1}\}  \\
&\qquad \times \int_{\mathbb{R}^n \setminus B_{R}(x_0)} \left(\frac{Z^{p_{+}(B_{\widetilde{R}}(x_0)\times B_{\widetilde{R}}(x_0)^c)-1}}{|y-x_0|^{n+sp_{+}(B_{\widetilde{R}}(x_0)\times B_{\widetilde{R}}(x_0)^c)}}  + \frac{Z^{p_{-}(B_{\widetilde{R}}(x_0)\times B_{\widetilde{R}}(x_0)^c)-1}}{|y-x_0|^{n+sp_{-}(B_{\widetilde{R}}(x_0)\times B_{\widetilde{R}}(x_0)^c)}} \right)\,\mathrm{d}y \\
& \leq C_0R^{-sp_{+}(B_{\widetilde{R}}(x_0)\times B_{\widetilde{R}}(x_0)^c)}Z^{p_{+}(B_{\widetilde{R}}(x_0)\times B_{\widetilde{R}}(x_0)^c)-1} \\
& \quad + C_0R^{-sp_{-}(B_{\widetilde{R}}(x_0)\times B_{\widetilde{R}}(x_0)^c)}Z^{p_{-}(B_{\widetilde{R}}(x_0)\times B_{\widetilde{R}}(x_0)^c)-1} \\
& =C_0(4^{j}\widetilde{R})^{-sp_{+}(B_{\widetilde{R}}(x_0)\times B_{\widetilde{R}}(x_0)^c)}(4^{\alpha j}2H)^{p_{+}(B_{\widetilde{R}}(x_0)\times B_{\widetilde{R}}(x_0)^c)-1} \\
& \quad + C_0(4^{j}\widetilde{R})^{-sp_{-}(B_{\widetilde{R}}(x_0)\times B_{\widetilde{R}}(x_0)^c)}(4^{\alpha j}2H)^{p_{-}(B_{\widetilde{R}}(x_0)\times B_{\widetilde{R}}(x_0)^c)-1} \\
& \leq C_04^{\frac{-sp_{+}(B_{\widetilde{R}}(x_0)\times B_{\widetilde{R}}(x_0)^c)}{2}j_0}H^{p_{+}(B_{\widetilde{R}}(x_0)\times B_{\widetilde{R}}(x_0)^c)-1}\widetilde{R}^{-sp_{+}(B_{\widetilde{R}}(x_0)\times B_{\widetilde{R}}(x_0)^c)} \\
& \quad +C_04^{\frac{-sp_{-}(B_{\widetilde{R}}(x_0)\times B_{\widetilde{R}}(x_0)^c)}{2}j_0}H^{p_{-}(B_{\widetilde{R}}(x_0)\times B_{\widetilde{R}}(x_0)^c)-1}\widetilde{R}^{-sp_{-}(B_{\widetilde{R}}(x_0)\times B_{\widetilde{R}}(x_0)^c)} \\
& \leq \frac{1}{2}(\delta H)^{p_+(B_{\widetilde{R}}(x_0)\times B_{\widetilde{R}}(x_0)^c) - 1} \widetilde{R}^{-sp_+(B_{\widetilde{R}}(x_0)\times B_{\widetilde{R}}(x_0)^c)}\\
&\quad + \frac{1}{2}(\delta H)^{p_-(B_{\widetilde{R}}(x_0)\times B_{\widetilde{R}}(x_0)^c)-1} \widetilde{R}^{-sp_-(B_{\widetilde{R}}(x_0)\times B_{\widetilde{R}}(x_0)^c)} \\
& \leq \frac{1}{2}(\delta H)^{p_+(B_{\widetilde{R}}(x_0)\times B_{\widetilde{R}}(x_0)) - 1} \widetilde{R}^{-sp_+(B_{\widetilde{R}}(x_0)\times B_{\widetilde{R}}(x_0))} \\
&\quad +\frac{1}{2}(\delta H)^{p_-(B_{\widetilde{R}}(x_0)\times B_{\widetilde{R}}(x_0))-1} \widetilde{R}^{-sp_-(B_{\widetilde{R}}(x_0)\times B_{\widetilde{R}}(x_0))},
\end{aligned}
\end{equation}
where we used the definition of $Z$, \eqref{eq:P2}, \eqref{eq:alphabounds} and \eqref{eq:j0}. 
Note that the constant $C_0$ comes from \eqref{eq:j0}. Combining the estimates of $J_1$ and $J_2$, proves \eqref{eq:assumption_tail2}.

Hence, we can apply \Cref{lem:growth}, which leads to
\[ u\geq m_j+\delta H = m_j+\delta \frac{M_{j}-m_j}{2} = m_j + \frac{\delta 4^{-\alpha j}Z}{2} > m_j+4^{-\alpha j}(1-4^{-\alpha})Z \quad \text{in } B_{\widetilde{R}/4}(x_0), \]
where we used \eqref{eq:alphabounds} in the last inequality. 
Hence, choosing $M_{j+1} = M_j$ and $m_{j+1}=m_j+4^{-\alpha j}(1-4^{-\alpha})Z$ proves \eqref{eq:sequencesMjmj} for the case \eqref{eq:Hoeldcase1}.

In the second case
\[ \left| B_{\frac{4^{-j}R}{2}}(x_0) \cap \left\{ u\geq m_{j} + \frac{M_j-m_j}{2} \right\} \right| < \frac{1}{2}\left| B_{\frac{4^{-j}R}{2}}(x_0) \right|, \]
we can proceed similarly and consider the function $v:=M_j-u$. In this case, we can choose the members of the sequences to be of the form $M_{j+1} = M_j-4^{-\alpha j}(1-4^{-\alpha})Z$ and $m_{j+1}=m_j$. This completes the construction of the sequences $(M_j)$ and $(m_j)$ and completes the proof of \eqref{eq:sequencesMjmj}. Now the local Hölder regularity follows in a standard way.
\end{proof}

\end{document}
